\section{Conclusion}
\label{sec:conclusion}

We evaluated a MERT-based AI-music detector under a protocol that varies the
generator and the corpus of real music independently, and found that these two
axes behave completely differently. Withholding an entire generator while
holding the real class fixed costs almost nothing: $F_1$ degrades by 0.037 and
0.023 across the two Suno/Udio directions, against 0.361 reported previously.
Replacing the corpus of real music collapses the same detector to near chance
and causes it to label every genuine recording as synthetic.

The reconciling explanation is that the detector encodes what real music looks
like in its training corpus rather than what generation leaves behind. Strong
cross-generator numbers are therefore not evidence that a detector has learned
to recognise synthesis, and a detector validated on cross-generator splits alone
may fail in deployment in the most costly direction available to it.

We also found that the conventional headline metric conceals this. Our
cross-corpus $F_1$ of 0.750 exceeds the published baseline of 0.629 and is
identical to the score of a classifier that labels every input synthetic.
Reporting the all-positive baseline and the false-positive rate alongside $F_1$
would have made the failure visible at a glance, and we recommend both as
routine practice.

Finally, the explanation layer produced two negative results worth carrying
forward. Gradient-based temporal explanation proved faithful for real audio and
unfaithful for generated audio, indicating that the evidence for synthesis is
not localised in time and that frequency-resolved explanation is the appropriate
tool. And the model's own layer-attention weights---the explanation an
attention-pooling architecture supplies for free---are flat, and their ordering
is negatively rank-correlated with exact Shapley attribution over the same
layers in every song we audited. Architectural attention should be reported as a
property of the architecture, not as an account of the decision.

The immediate next step is an evaluation in which two architecturally distinct
generator families are paired with a \emph{shared} corpus of real music, which
would separate the corpus effect from the generator effect that this work shows
are conflated in current practice.
